
@@M0@@ sur @@M1@@. Cet
homomorphisme définit une immersion fermée @@M2@@, régulière d'après 1.8 ii).

D'après 1.8 iii), @@M3@@ est parfait en tant que @@M4@@-Module gradué. En particulier, @@M5@@ est de présentation finie comme Module gradué. Soit donc:
@@M6@@
une présentation finie graduée de @@M7@@ sur @@M8@@. On obtient sur @@M9@@ la suite exacte:
@@M10@@
où @@M11@@, @@M12@@, et sont donc localement libres de type fini sur @@M13@@. On peut donc écrire pour tout @@M14@@ le diagramme commutatif:
@@M15@@
où @@M16@@ est le morphisme structural: @@M17@@, et @@M18@@, @@M19@@ les composantes de degré @@M20@@ de @@M21@@, @@M22@@. La ligne supérieure est exacte par hypothèse; les deux premières flèches verticales sont des isomorphismes à cause des résultats classiques pour les images directes dans @@M23@@, @@M24@@ étant un fibré projectif (EGA III, 2.1.15); enfin, la dernière ligne est exacte si @@M25@@ est assez grand, car le noyau de @@M26@@ est un @@M27@@-Module parfait, et on peut appliquer le lemme rappelé plus haut, qui montre que ce noyau, tensorisé par @@M28@@ pour @@M29@@ assez grand, est @@M30@@-acyclique; de même pour le noyau de @@M31@@. Il en résulte alors que pour @@M32@@ assez grand, @@M33@@ est un isomorphisme. Or @@M34@@; on a
